\documentclass[12pt]{article}
\usepackage{amsmath,amssymb}
\usepackage{geometry}
\usepackage{enumitem}
\geometry{margin=0.8in}

\title{Algebra II: Vectors and Parametric Equations}
\author{}
\date{}

\begin{document}
\maketitle

\section*{Formula and Concept Sheet}

\subsection*{Vectors in Two and Three Dimensions}
A vector records direction and length. Its components show the change in each coordinate:
\[
  \mathbf{v}=\langle a,b\rangle \quad\text{in 2D},
  \qquad
  \mathbf{v}=\langle a,b,c\rangle \quad\text{in 3D}.
\]
If $P=(x_1,y_1)$ and $Q=(x_2,y_2)$, then
\[
  \overrightarrow{PQ}=\langle x_2-x_1,y_2-y_1\rangle.
\]
In 3D, include the change $z_2-z_1$ as the third component.

For vectors with matching dimensions,
\[
\begin{aligned}
  \langle a,b\rangle+\langle c,d\rangle&=\langle a+c,b+d\rangle,\\
  k\langle a,b\rangle&=\langle ka,kb\rangle.
\end{aligned}
\]
The same rules apply component by component in 3D.

\subsection*{Length, Dot Product, and Angle}
The length, or magnitude, of a vector is
\[
  \|\langle a,b\rangle\|=\sqrt{a^2+b^2},
  \qquad
  \|\langle a,b,c\rangle\|=\sqrt{a^2+b^2+c^2}.
\]
The dot product is
\[
  \langle a,b\rangle\cdot\langle c,d\rangle=ac+bd
\]
and, in 3D,
\[
  \langle a,b,c\rangle\cdot\langle d,e,f\rangle=ad+be+cf.
\]
For nonzero vectors $\mathbf{u}$ and $\mathbf{v}$, the angle $\phi$ between them satisfies
\[
  \cos\phi=\frac{\mathbf{u}\cdot\mathbf{v}}
  {\|\mathbf{u}\|\,\|\mathbf{v}\|},
  \qquad 0\le\phi\le\pi.
\]
Useful tests:
\[
  \mathbf{u}\cdot\mathbf{v}=0 \Longrightarrow \mathbf{u}\perp\mathbf{v},
  \qquad
  \mathbf{u}=k\mathbf{v} \Longrightarrow \mathbf{u}\parallel\mathbf{v}.
\]

\subsection*{Components Parametrized by $t$}
A vector-valued rule lets every component depend on one parameter:
\[
  \mathbf{r}(t)=\langle x(t),y(t)\rangle
  \quad\text{or}\quad
  \mathbf{r}(t)=\langle x(t),y(t),z(t)\rangle.
\]
Each value of $t$ gives one point. The parameter need not represent time unless the context says it does.

A line through $P=(x_0,y_0)$ with direction vector $\langle a,b\rangle$ can be parametrized by
\[
  \mathbf{r}(t)=\langle x_0,y_0\rangle+t\langle a,b\rangle,
\]
or equivalently,
\[
  x=x_0+at,\qquad y=y_0+bt.
\]
To parametrize the line through points $P$ and $Q$, use the direction vector $\overrightarrow{PQ}$. Different nonzero scalar multiples of a direction vector describe the same line.

In 3D, a line through $(x_0,y_0,z_0)$ with direction $\langle a,b,c\rangle$ is
\[
  x=x_0+at,\qquad y=y_0+bt,\qquad z=z_0+ct.
\]
To remove a parameter, solve one component equation for $t$ and substitute into the others.

\subsection*{Components Parametrized by $\theta$}
A circle of radius $R>0$ centered at $(h,k)$ has the standard parametrization
\[
  x=h+R\cos\theta,
  \qquad
  y=k+R\sin\theta,
  \qquad 0\le\theta<2\pi.
\]
As $\theta$ increases from $0$, this starts at $(h+R,k)$ and travels counterclockwise. Every point satisfies
\[
  (x-h)^2+(y-k)^2=R^2
\]
because $\cos^2\theta+\sin^2\theta=1$. Swapping sine and cosine, inserting negative signs, or shifting $\theta$ can change the starting point or direction without changing the circle.

\clearpage
\section*{30-Minute Multiple-Choice Practice}
\textbf{Instructions.} Complete all 12 questions in 30 minutes. Select the best answer for each question. A calculator may be used, but give exact values unless a decimal is requested.

\begin{enumerate}[leftmargin=*,itemsep=1.1em]
  \item If $\mathbf{u}(t)=\langle 2t-1,t^2-4\rangle$, what is $\|\mathbf{u}(3)\|$?
  \begin{enumerate}[label=\Alph*.,itemsep=0pt]
    \item $5$ \item $5\sqrt2$ \item $10$ \item $10\sqrt2$
  \end{enumerate}

  \item Which parametrization describes the line through $(-2,5)$ and $(4,-1)$?
  \begin{enumerate}[label=\Alph*.,itemsep=0pt]
    \item $x=-2+t,\ y=5-t$
    \item $x=-2+t,\ y=5+t$
    \item $x=4-2t,\ y=-1+5t$
    \item $x=-2+4t,\ y=5-t$
  \end{enumerate}

  \item The line $x=3-2t$, $y=-1+4t$ intersects the line $y=x$. What is the point of intersection?
  \begin{enumerate}[label=\Alph*.,itemsep=0pt]
    \item $\left(\dfrac13,\dfrac13\right)$
    \item $\left(\dfrac53,\dfrac53\right)$
    \item $(1,1)$
    \item $(3,3)$
  \end{enumerate}

  \item Which Cartesian equation is traced by $x=1+3t$, $y=-2+6t$?
  \begin{enumerate}[label=\Alph*.,itemsep=0pt]
    \item $y=2x-4$ \item $y=2x-2$
    \item $y=3x-5$ \item $y=6x-8$
  \end{enumerate}

  \item Which parametrization traces the circle centered at $(-3,2)$ with radius $5$ counterclockwise, starting at its rightmost point when $\theta=0$?
  \begin{enumerate}[label=\Alph*.,itemsep=0pt]
    \item $x=-3+5\cos\theta,\ y=2+5\sin\theta$
    \item $x=-3+5\sin\theta,\ y=2+5\cos\theta$
    \item $x=3+5\cos\theta,\ y=-2+5\sin\theta$
    \item $x=-3-5\cos\theta,\ y=2+5\sin\theta$
  \end{enumerate}

  \item The equations $x=4+3\sin\theta$, $y=-1-3\cos\theta$ trace which motion as $\theta$ increases from $0$?
  \begin{enumerate}[label=\Alph*.,itemsep=0pt]
    \item A radius-3 circle centered at $(4,-1)$, starting at the bottom and moving counterclockwise
    \item A radius-3 circle centered at $(4,-1)$, starting at the top and moving clockwise
    \item A radius-3 circle centered at $(-4,1)$, starting at the bottom and moving counterclockwise
    \item A radius-9 circle centered at $(4,-1)$, starting at the right and moving clockwise
  \end{enumerate}

  \item What is the angle between $\langle3,4\rangle$ and $\langle-4,3\rangle$?
  \begin{enumerate}[label=\Alph*.,itemsep=0pt]
    \item $0^\circ$ \item $45^\circ$ \item $90^\circ$ \item $180^\circ$
  \end{enumerate}

  \item What is the angle between $\langle1,\sqrt3\rangle$ and $\langle2,0\rangle$?
  \begin{enumerate}[label=\Alph*.,itemsep=0pt]
    \item $30^\circ$ \item $45^\circ$ \item $60^\circ$ \item $120^\circ$
  \end{enumerate}

  \item For what value of $k$ are $\langle k,2\rangle$ and $\langle3,-6\rangle$ perpendicular?
  \begin{enumerate}[label=\Alph*.,itemsep=0pt]
    \item $-4$ \item $-2$ \item $2$ \item $4$
  \end{enumerate}

  \item What is the distance from $P=(1,-2,3)$ to $Q=(5,1,-3)$?
  \begin{enumerate}[label=\Alph*.,itemsep=0pt]
    \item $\sqrt{29}$ \item $\sqrt{52}$ \item $\sqrt{61}$ \item $13$
  \end{enumerate}

  \item What is the angle between $\langle1,2,2\rangle$ and $\langle2,1,-2\rangle$?
  \begin{enumerate}[label=\Alph*.,itemsep=0pt]
    \item $30^\circ$ \item $60^\circ$ \item $90^\circ$ \item $120^\circ$
  \end{enumerate}

  \item A line passes through $(2,-1,4)$ and $(-1,5,1)$. At what point does the line intersect the $xy$-plane?
  \begin{enumerate}[label=\Alph*.,itemsep=0pt]
    \item $(-2,7,0)$ \item $(-1,5,0)$
    \item $(0,3,0)$ \item $(2,-1,0)$
  \end{enumerate}
\end{enumerate}

\end{document}
